\section{Space Complexity --- Machines --- ARMWidth Example}
%%% generated by the blueprint pipeline from TCSlib.Complexity.SpaceComplexity.Machines.ARMWidthExample
\subsection*{Space Complexity --- Machines --- ARMWidth Example: Introduction}

Declarations of the Lean module \texttt{TCSlib.Complexity.SpaceComplexity.Machines.ARMWidthExample}.

\subsection*{Space Complexity --- Machines --- ARMWidth Example: Results}

\begin{definition}[Bounded search language]
\lean{Complexity.LogProg.WidthExample.searchLang}
\label{Complexity.LogProg.WidthExample.searchLang}
\leanok
\uses{Turing.pairEncode}
Let $P$ be a language of Boolean words. The bounded search language consists of those
Boolean words $x$ for which there is a natural number $u<2^{|x|}$ such that the pair
encoding of $x$ and the binary word representing $u$ belongs to $P$.
\end{definition}

\begin{definition}[Finite set of control states]
\lean{Complexity.LogProg.WidthExample.Lb}
\label{Complexity.LogProg.WidthExample.Lb}
\leanok
Introduce a finite set of seventeen distinct control states, named init, scan, incI,
dbl, decI, d1, d1a, d1b, d1c, d2, d2a, d2b, enum, chk, nxt, acc, and rej.
\end{definition}

\begin{definition}[Four-register search program]
\lean{Complexity.LogProg.WidthExample.searchA}
\label{Complexity.LogProg.WidthExample.searchA}
\leanok
\uses{Complexity.LogProg.ARM, Complexity.LogProg.WidthExample.Lb}
Consider a four-register program with two oracle deciders and the control states listed
for the search machine. It starts at init, increments register 1 and enters scan; scan
calls decider 0 on the whole-input virtual query formed from register 0, proceeding to
incI on yes and dbl on no; incI increments register 0 and returns to scan. At dbl it
proceeds to enum if register 0 is zero, and otherwise decrements register 0 and enters
decI; decI decrements register 1 and enters d1. In d1, if register 1 is zero it enters
d2, and otherwise it decrements register 1, increments register 2 twice through d1a,
d1b, and d1c, and returns to d1. In d2, if register 2 is zero it returns to dbl, and
otherwise it decrements register 2, increments register 1, and returns to d2. At enum it
rejects if registers 3 and 1 are equal, and otherwise enters chk; chk calls decider 1 on
the whole-input virtual query formed from register 3, accepting on yes and entering nxt
on no; nxt increments register 3 and returns to enum. The accepting and rejecting states
halt with answers true and false, respectively.
\end{definition}

\begin{definition}[Two-decider language family]
\lean{Complexity.LogProg.WidthExample.decs}
\label{Complexity.LogProg.WidthExample.decs}
\leanok
\uses{Complexity.inputLenLang}
For a language $P$ of Boolean words, define a family of two languages indexed by
$j\in\{0,1\}$: the language at index $0$ is the input-length language, and the language
at index $1$ is $P$.
\end{definition}

\begin{definition}[Oracle for the two languages]
\lean{Complexity.LogProg.WidthExample.orc}
\label{Complexity.LogProg.WidthExample.orc}
\leanok
\uses{Complexity.LogProg.WidthExample.decs, Turing.MultiTapeTM.indicator}
For a language $P$ of Boolean words, define an oracle with two indices that answers true
on a Boolean-word query exactly when that query belongs to the language at that index in
the family consisting of the input-length language and $P$.
\end{definition}

\begin{definition}[Four register values]
\lean{Complexity.LogProg.WidthExample.rf}
\label{Complexity.LogProg.WidthExample.rf}
\leanok
For natural numbers $i,p,t,u$, define a four-entry register assignment whose values, in
order, are $i,p,t,u$.
\end{definition}

\begin{lemma}[Updating the First Coordinate]
\lean{Complexity.LogProg.WidthExample.rf_upd0}
\label{Complexity.LogProg.WidthExample.rf_upd0}
\leanok
\uses{Complexity.LogProg.WidthExample.rf}
\difficulty{4}
For natural numbers $i,p,t,u,a$, let $r(i,p,t,u)$ be the function on four indices with
values $i,p,t,u$ in order. Replacing its value at index $0$ by $a$ gives $r(a,p,t,u)$.
\end{lemma}

\begin{lemma}[Updating the Second Coordinate]
\lean{Complexity.LogProg.WidthExample.rf_upd1}
\label{Complexity.LogProg.WidthExample.rf_upd1}
\leanok
\uses{Complexity.LogProg.WidthExample.rf}
\difficulty{4}
For natural numbers $i,p,t,u,a$, let $r(i,p,t,u)$ be the function on four indices with
values $i,p,t,u$ in order. Replacing its value at index $1$ by $a$ gives $r(i,a,t,u)$.
\end{lemma}

\begin{lemma}[Updating the Third Coordinate]
\lean{Complexity.LogProg.WidthExample.rf_upd2}
\label{Complexity.LogProg.WidthExample.rf_upd2}
\leanok
\uses{Complexity.LogProg.WidthExample.rf}
\difficulty{4}
For natural numbers $i,p,t,u,a$, let $r(i,p,t,u)$ be the function on four indices with
values $i,p,t,u$ in order. Replacing its value at index $2$ by $a$ gives $r(i,p,a,u)$.
\end{lemma}

\begin{lemma}[Updating the Fourth Coordinate]
\lean{Complexity.LogProg.WidthExample.rf_upd3}
\label{Complexity.LogProg.WidthExample.rf_upd3}
\leanok
\uses{Complexity.LogProg.WidthExample.rf}
\difficulty{4}
For natural numbers $i,p,t,u,a$, let $r(i,p,t,u)$ be the function on four indices with
values $i,p,t,u$ in order. Replacing its value at index $3$ by $a$ gives $r(i,p,t,a)$.
\end{lemma}

\begin{lemma}[First Register Coordinate]
\lean{Complexity.LogProg.WidthExample.rf_0}
\label{Complexity.LogProg.WidthExample.rf_0}
\leanok
\uses{Complexity.LogProg.WidthExample.rf}
\difficulty{1}
For all natural numbers $i,p,t,u$, the first register of the four-register tuple
$(i,p,t,u)$ is $i$.
\end{lemma}

\begin{lemma}[Second Register Coordinate]
\lean{Complexity.LogProg.WidthExample.rf_1}
\label{Complexity.LogProg.WidthExample.rf_1}
\leanok
\uses{Complexity.LogProg.WidthExample.rf}
\difficulty{1}
For all natural numbers $i,p,t,u$, the second register of the four-register tuple
$(i,p,t,u)$ is $p$.
\end{lemma}

\begin{lemma}[Third Register Coordinate]
\lean{Complexity.LogProg.WidthExample.rf_2}
\label{Complexity.LogProg.WidthExample.rf_2}
\leanok
\uses{Complexity.LogProg.WidthExample.rf}
\difficulty{1}
For all natural numbers $i,p,t,u$, the third register of the four-register tuple
$(i,p,t,u)$ is $t$.
\end{lemma}

\begin{lemma}[Fourth Register Coordinate]
\lean{Complexity.LogProg.WidthExample.rf_3}
\label{Complexity.LogProg.WidthExample.rf_3}
\leanok
\uses{Complexity.LogProg.WidthExample.rf}
\difficulty{1}
For all natural numbers $i,p,t,u$, the fourth register of the four-register tuple
$(i,p,t,u)$ is $u$.
\end{lemma}

\begin{definition}[Bounded pre-search configurations]
\lean{Complexity.LogProg.WidthExample.G}
\label{Complexity.LogProg.WidthExample.G}
\leanok
\uses{Complexity.LogProg.AConf, Complexity.LogProg.PreS, Complexity.LogProg.WidthExample.Lb, Complexity.LogProg.WidthExample.searchA}
For a Boolean-word input $x$ and a natural-number bound $B$, let $G(x,B)$ be the set of
abstract configurations $a$ for which $a$ satisfies the precondition for the search
program on input $x$ and every register value in $a$ is at most $B$.
\end{definition}

\begin{lemma}[Register Configuration Invariant]
\lean{Complexity.LogProg.WidthExample.guard_rf}
\label{Complexity.LogProg.WidthExample.guard_rf}
\leanok
\uses{Complexity.LogProg.PreS, Complexity.LogProg.WidthExample.G, Complexity.LogProg.WidthExample.Lb, Complexity.LogProg.WidthExample.rf, Complexity.LogProg.WidthExample.searchA}
\difficulty{4}
For any input word $x$, bound $B$, program location $l$, and natural numbers $i,p,t,u$
with $i,p,t,u\le B$, the configuration at location $l$ with registers $(i,p,t,u)$ and no
answer satisfies the search program's precondition and has every register at most $B$.
\end{lemma}

\begin{lemma}[Register values during a decrement loop]
\lean{Complexity.LogProg.WidthExample.d1_loop}
\label{Complexity.LogProg.WidthExample.d1_loop}
\leanok
\uses{Complexity.LogProg.AReach, Complexity.LogProg.AReach.refl, Complexity.LogProg.AReach.step, Complexity.LogProg.WidthExample.G, Complexity.LogProg.WidthExample.guard_rf, Complexity.LogProg.WidthExample.orc, Complexity.LogProg.WidthExample.rf, Complexity.LogProg.WidthExample.rf_1, Complexity.LogProg.WidthExample.rf_2, Complexity.LogProg.WidthExample.rf_upd1, Complexity.LogProg.WidthExample.rf_upd2, Complexity.LogProg.WidthExample.searchA, Complexity.LogProg.astep}
\difficulty{5}
Let $P$ be a language of bit strings, let $x$ be a bit string, and let $B,i,u\in\mathbb
N$ satisfy $i\le B$ and $u\le B$. For every $a,b\in\mathbb N$ with $b+2a\le B$, the
program starts at its first decrement-loop label with register values $(i,a,b,u)$ and,
after finitely many steps, reaches the next loop label with values $(i,0,b+2a,u)$ and no
answer recorded. At every step before reaching that configuration, all four register
values are at most $B$.
\end{lemma}

\begin{lemma}[D2 Loop Register Update]
\lean{Complexity.LogProg.WidthExample.d2_loop}
\label{Complexity.LogProg.WidthExample.d2_loop}
\leanok
\uses{Complexity.LogProg.AReach, Complexity.LogProg.AReach.refl, Complexity.LogProg.AReach.step, Complexity.LogProg.WidthExample.G, Complexity.LogProg.WidthExample.guard_rf, Complexity.LogProg.WidthExample.orc, Complexity.LogProg.WidthExample.rf, Complexity.LogProg.WidthExample.rf_1, Complexity.LogProg.WidthExample.rf_2, Complexity.LogProg.WidthExample.rf_upd1, Complexity.LogProg.WidthExample.rf_upd2, Complexity.LogProg.WidthExample.searchA, Complexity.LogProg.astep}
\difficulty{4}
For any language $P$, input word $x$, and bound $B$, suppose $i,u\le B$. For every
$b,p\in\mathbb{N}$ with $p+b\le B$, the search program, using its oracle determined by
$P$, can reach the configuration at location $\mathrm{dbl}$ with registers $(i,p+b,0,u)$
from location $\mathrm{d2}$ with registers $(i,p,b,u)$, without leaving the
configurations satisfying the program precondition and the bound $B$ on every register
at any earlier step.
\end{lemma}

\begin{lemma}[Doubling Loop Register Update]
\lean{Complexity.LogProg.WidthExample.dbl_loop}
\label{Complexity.LogProg.WidthExample.dbl_loop}
\leanok
\uses{Complexity.LogProg.AReach, Complexity.LogProg.AReach.refl, Complexity.LogProg.AReach.step, Complexity.LogProg.AReach.trans, Complexity.LogProg.WidthExample.G, Complexity.LogProg.WidthExample.d1_loop, Complexity.LogProg.WidthExample.d2_loop, Complexity.LogProg.WidthExample.guard_rf, Complexity.LogProg.WidthExample.orc, Complexity.LogProg.WidthExample.rf, Complexity.LogProg.WidthExample.rf_0, Complexity.LogProg.WidthExample.rf_upd0, Complexity.LogProg.WidthExample.searchA, Complexity.LogProg.astep}
\difficulty{5}
For any language $P$, input word $x$, and natural number $n$ with $2^n\le B$, whenever
$a+k=n$, the search program with oracle determined by $P$ can reach location
$\mathrm{enum}$ with registers $(0,2^n,0,0)$ from location $\mathrm{dbl}$ with registers
$(k,2^a,0,0)$. Throughout the run before reaching that configuration, the program
precondition holds and every register is at most $B$.
\end{lemma}

\begin{lemma}[Enumeration Loop Halting]
\lean{Complexity.LogProg.WidthExample.enum_loop}
\label{Complexity.LogProg.WidthExample.enum_loop}
\leanok
\uses{Complexity.LogProg.AHalt, Complexity.LogProg.AHalt.ret, Complexity.LogProg.AHalt.step, Complexity.LogProg.WidthExample.G, Complexity.LogProg.WidthExample.decs, Complexity.LogProg.WidthExample.guard_rf, Complexity.LogProg.WidthExample.orc, Complexity.LogProg.WidthExample.rf, Complexity.LogProg.WidthExample.rf_1, Complexity.LogProg.WidthExample.rf_3, Complexity.LogProg.WidthExample.rf_upd3, Complexity.LogProg.WidthExample.searchA, Complexity.LogProg.astep, Complexity.LogProg.indicator_eq_decide, Complexity.LogProg.vword_whole₁, Turing.pairEncode}
\difficulty{6}
For any language $P$, input word $x$, and natural number $n$ with $2^n\le B$, let $k,u$
satisfy $u+k=2^n$. For either Boolean value $b$ such that $b$ is true exactly when there
exists $w$ with $u\le w<2^n$ and the encoded pair of $x$ and the binary representation
of $w$ belongs to $P$, the search program with oracle determined by $P$ halts with
answer $b$, while satisfying the program precondition and keeping every register at most
$B$ at every earlier step.
\end{lemma}

\begin{theorem}[Search Preserves Polynomial Space]
\lean{Complexity.LogProg.WidthExample.searchLang_mem_PSPACE}
\label{Complexity.LogProg.WidthExample.searchLang_mem_PSPACE}
\leanok
\uses{Complexity.LOGSPACE_subset_PSPACE, Complexity.LogProg.AHalt.step, Complexity.LogProg.AReach.halt, Complexity.LogProg.WidthExample.G, Complexity.LogProg.WidthExample.dbl_loop, Complexity.LogProg.WidthExample.decs, Complexity.LogProg.WidthExample.enum_loop, Complexity.LogProg.WidthExample.guard_rf, Complexity.LogProg.WidthExample.orc, Complexity.LogProg.WidthExample.rf, Complexity.LogProg.WidthExample.rf_1, Complexity.LogProg.WidthExample.rf_upd0, Complexity.LogProg.WidthExample.rf_upd1, Complexity.LogProg.WidthExample.searchA, Complexity.LogProg.WidthExample.searchLang, Complexity.LogProg.arm_decides_pspace_val, Complexity.LogProg.astep, Complexity.LogProg.indicator_eq_decide, Complexity.LogProg.inputLength_loop, Complexity.PSPACE, Complexity.indicator_indexLang_pairEncode, Complexity.inputLenLang, Complexity.inputLenLang_mem, Turing.pairEncode}
\difficulty{5}
For a language $P$ of Boolean words, define its search language to consist of those
Boolean words $x$ for which there is a natural number $u<2^{|x|}$ such that the encoded
pair of $x$ and the binary expansion of $u$ belongs to $P$. If $P$ is decidable by a
finite-state multi-tape Turing machine using at most a constant multiple of $n^c+1$
work-tape cells on inputs of length $n$, for some natural number $c$, then the search
language of $P$ also has such a polynomial-space decision procedure.
\end{theorem}
